\providecommand{\eventanchor}[2]{\hypertarget{evt:#1}{#2}}
\section{}

\paragraph{983年}\eventanchor{NS-0983-REGNAL-YEAR}{NS-0983-REGNAL-YEAR}　宋以宋太宗在位第7年作为诏令、赋役与外交文书的纪年框架。账本列出《宋史》卷1—23〈本纪〉；《续资治通鉴长编》本年条；《宋会要辑稿》相关门类作为来源线索；该材料可核对时间与制度称谓，但有关执行范围、群体经验、军数或后果，仍须与地方文书、物质遗存及相关政权记录互证。

\paragraph{984年}\eventanchor{NS-0984-REGNAL-YEAR}{NS-0984-REGNAL-YEAR}　宋以宋太宗在位第8年作为诏令、赋役与外交文书的纪年框架。账本列出《宋史》卷1—23〈本纪〉；《续资治通鉴长编》本年条；《宋会要辑稿》相关门类作为来源线索；该材料可核对时间与制度称谓，但有关执行范围、群体经验、军数或后果，仍须与地方文书、物质遗存及相关政权记录互证。

\paragraph{985年}\eventanchor{NS-0985-REGNAL-YEAR}{NS-0985-REGNAL-YEAR}　宋以宋太宗在位第9年作为诏令、赋役与外交文书的纪年框架。账本列出《宋史》卷1—23〈本纪〉；《续资治通鉴长编》本年条；《宋会要辑稿》相关门类作为来源线索；该材料可核对时间与制度称谓，但有关执行范围、群体经验、军数或后果，仍须与地方文书、物质遗存及相关政权记录互证。

\paragraph{986年}\eventanchor{NS-0986-REGNAL-YEAR}{NS-0986-REGNAL-YEAR}　宋以宋太宗在位第10年作为诏令、赋役与外交文书的纪年框架。账本列出《宋史》卷1—23〈本纪〉；《续资治通鉴长编》本年条；《宋会要辑稿》相关门类作为来源线索；该材料可核对时间与制度称谓，但有关执行范围、群体经验、军数或后果，仍须与地方文书、物质遗存及相关政权记录互证。

\paragraph{986年}\eventanchor{NS-0986-YONGXI-CAMPAIGN}{NS-0986-YONGXI-CAMPAIGN}　宋以三路军北伐辽，随后遭挫撤军。账本列出《宋史》卷4〈太宗本纪二〉；《辽史》〈圣宗本纪〉；《续资治通鉴长编》雍熙三年条作为来源线索；该材料可核对时间与制度称谓，但有关执行范围、群体经验、军数或后果，仍须与地方文书、物质遗存及相关政权记录互证。

\paragraph{987年}\eventanchor{NS-0987-REGNAL-YEAR}{NS-0987-REGNAL-YEAR}　宋以宋太宗在位第11年作为诏令、赋役与外交文书的纪年框架。账本列出《宋史》卷1—23〈本纪〉；《续资治通鉴长编》本年条；《宋会要辑稿》相关门类作为来源线索；该材料可核对时间与制度称谓，但有关执行范围、群体经验、军数或后果，仍须与地方文书、物质遗存及相关政权记录互证。

\paragraph{988年}\eventanchor{NS-0988-REGNAL-YEAR}{NS-0988-REGNAL-YEAR}　宋以宋太宗在位第12年作为诏令、赋役与外交文书的纪年框架。账本列出《宋史》卷1—23〈本纪〉；《续资治通鉴长编》本年条；《宋会要辑稿》相关门类作为来源线索；该材料可核对时间与制度称谓，但有关执行范围、群体经验、军数或后果，仍须与地方文书、物质遗存及相关政权记录互证。

\paragraph{989年}\eventanchor{NS-0989-REGNAL-YEAR}{NS-0989-REGNAL-YEAR}　宋以宋太宗在位第13年作为诏令、赋役与外交文书的纪年框架。账本列出《宋史》卷1—23〈本纪〉；《续资治通鉴长编》本年条；《宋会要辑稿》相关门类作为来源线索；该材料可核对时间与制度称谓，但有关执行范围、群体经验、军数或后果，仍须与地方文书、物质遗存及相关政权记录互证。

\paragraph{990年}\eventanchor{NS-0990-REGNAL-YEAR}{NS-0990-REGNAL-YEAR}　宋以宋太宗在位第14年作为诏令、赋役与外交文书的纪年框架。账本列出《宋史》卷1—23〈本纪〉；《续资治通鉴长编》本年条；《宋会要辑稿》相关门类作为来源线索；该材料可核对时间与制度称谓，但有关执行范围、群体经验、军数或后果，仍须与地方文书、物质遗存及相关政权记录互证。

\paragraph{991年}\eventanchor{NS-0991-REGNAL-YEAR}{NS-0991-REGNAL-YEAR}　宋以宋太宗在位第15年作为诏令、赋役与外交文书的纪年框架。账本列出《宋史》卷1—23〈本纪〉；《续资治通鉴长编》本年条；《宋会要辑稿》相关门类作为来源线索；该材料可核对时间与制度称谓，但有关执行范围、群体经验、军数或后果，仍须与地方文书、物质遗存及相关政权记录互证。

\paragraph{992年}\eventanchor{NS-0992-REGNAL-YEAR}{NS-0992-REGNAL-YEAR}　宋以宋太宗在位第16年作为诏令、赋役与外交文书的纪年框架。账本列出《宋史》卷1—23〈本纪〉；《续资治通鉴长编》本年条；《宋会要辑稿》相关门类作为来源线索；该材料可核对时间与制度称谓，但有关执行范围、群体经验、军数或后果，仍须与地方文书、物质遗存及相关政权记录互证。

\paragraph{993年}\eventanchor{NS-0993-REGNAL-YEAR}{NS-0993-REGNAL-YEAR}　宋以宋太宗在位第17年作为诏令、赋役与外交文书的纪年框架。账本列出《宋史》卷1—23〈本纪〉；《续资治通鉴长编》本年条；《宋会要辑稿》相关门类作为来源线索；该材料可核对时间与制度称谓，但有关执行范围、群体经验、军数或后果，仍须与地方文书、物质遗存及相关政权记录互证。

\paragraph{993年}\eventanchor{NS-0993-WANG-XIAOBO}{NS-0993-WANG-XIAOBO}　王小波在成都平原起事，提出均贫富诉求。账本列出《宋史》卷4〈太宗本纪二〉；《宋史》卷464〈王小波传〉；《续资治通鉴长编》淳化四年条作为来源线索；该材料可核对时间与制度称谓，但有关执行范围、群体经验、军数或后果，仍须与地方文书、物质遗存及相关政权记录互证。

\paragraph{994年}\eventanchor{NS-0994-LI-SHUN}{NS-0994-LI-SHUN}　李顺继领蜀地起事并一度占据成都，后被宋军平定。账本列出《宋史》卷4〈太宗本纪二〉；《宋史》卷464〈王小波传〉；《续资治通鉴长编》淳化五年条作为来源线索；该材料可核对时间与制度称谓，但有关执行范围、群体经验、军数或后果，仍须与地方文书、物质遗存及相关政权记录互证。

\paragraph{994年}\eventanchor{NS-0994-REGNAL-YEAR}{NS-0994-REGNAL-YEAR}　宋以宋太宗在位第18年作为诏令、赋役与外交文书的纪年框架。账本列出《宋史》卷1—23〈本纪〉；《续资治通鉴长编》本年条；《宋会要辑稿》相关门类作为来源线索；该材料可核对时间与制度称谓，但有关执行范围、群体经验、军数或后果，仍须与地方文书、物质遗存及相关政权记录互证。

\paragraph{995年}\eventanchor{NS-0995-REGNAL-YEAR}{NS-0995-REGNAL-YEAR}　宋以宋太宗在位第19年作为诏令、赋役与外交文书的纪年框架。账本列出《宋史》卷1—23〈本纪〉；《续资治通鉴长编》本年条；《宋会要辑稿》相关门类作为来源线索；该材料可核对时间与制度称谓，但有关执行范围、群体经验、军数或后果，仍须与地方文书、物质遗存及相关政权记录互证。

\paragraph{996年}\eventanchor{NS-0996-REGNAL-YEAR}{NS-0996-REGNAL-YEAR}　宋以宋太宗在位第20年作为诏令、赋役与外交文书的纪年框架。账本列出《宋史》卷1—23〈本纪〉；《续资治通鉴长编》本年条；《宋会要辑稿》相关门类作为来源线索；该材料可核对时间与制度称谓，但有关执行范围、群体经验、军数或后果，仍须与地方文书、物质遗存及相关政权记录互证。

\paragraph{997年}\eventanchor{NS-0997-REGNAL-YEAR}{NS-0997-REGNAL-YEAR}　宋以宋太宗在位第21年作为诏令、赋役与外交文书的纪年框架。账本列出《宋史》卷1—23〈本纪〉；《续资治通鉴长编》本年条；《宋会要辑稿》相关门类作为来源线索；该材料可核对时间与制度称谓，但有关执行范围、群体经验、军数或后果，仍须与地方文书、物质遗存及相关政权记录互证。

\paragraph{997年三月}\eventanchor{NS-0997-ZHENZONG-ACCESSION}{NS-0997-ZHENZONG-ACCESSION}　太宗去世，赵恒继位为真宗。账本列出《宋史》卷5〈真宗本纪一〉；《续资治通鉴长编》至道三年三月条；《宋会要辑稿·礼》即位条作为来源线索；该材料可核对时间与制度称谓，但有关执行范围、群体经验、军数或后果，仍须与地方文书、物质遗存及相关政权记录互证。

\paragraph{998年}\eventanchor{NS-0998-REGNAL-YEAR}{NS-0998-REGNAL-YEAR}　宋以宋真宗在位第1年作为诏令、赋役与外交文书的纪年框架。账本列出《宋史》卷1—23〈本纪〉；《续资治通鉴长编》本年条；《宋会要辑稿》相关门类作为来源线索；该材料可核对时间与制度称谓，但有关执行范围、群体经验、军数或后果，仍须与地方文书、物质遗存及相关政权记录互证。

\paragraph{999年}\eventanchor{NS-0999-REGNAL-YEAR}{NS-0999-REGNAL-YEAR}　宋以宋真宗在位第2年作为诏令、赋役与外交文书的纪年框架。账本列出《宋史》卷1—23〈本纪〉；《续资治通鉴长编》本年条；《宋会要辑稿》相关门类作为来源线索；该材料可核对时间与制度称谓，但有关执行范围、群体经验、军数或后果，仍须与地方文书、物质遗存及相关政权记录互证。

\paragraph{1000年}\eventanchor{NS-1000-REGNAL-YEAR}{NS-1000-REGNAL-YEAR}　宋以宋真宗在位第3年作为诏令、赋役与外交文书的纪年框架。账本列出《宋史》卷1—23〈本纪〉；《续资治通鉴长编》本年条；《宋会要辑稿》相关门类作为来源线索；该材料可核对时间与制度称谓，但有关执行范围、群体经验、军数或后果，仍须与地方文书、物质遗存及相关政权记录互证。

\paragraph{1001年}\eventanchor{NS-1001-REGNAL-YEAR}{NS-1001-REGNAL-YEAR}　宋以宋真宗在位第4年作为诏令、赋役与外交文书的纪年框架。账本列出《宋史》卷1—23〈本纪〉；《续资治通鉴长编》本年条；《宋会要辑稿》相关门类作为来源线索；该材料可核对时间与制度称谓，但有关执行范围、群体经验、军数或后果，仍须与地方文书、物质遗存及相关政权记录互证。

\paragraph{1002年}\eventanchor{NS-1002-REGNAL-YEAR}{NS-1002-REGNAL-YEAR}　宋以宋真宗在位第5年作为诏令、赋役与外交文书的纪年框架。账本列出《宋史》卷1—23〈本纪〉；《续资治通鉴长编》本年条；《宋会要辑稿》相关门类作为来源线索；该材料可核对时间与制度称谓，但有关执行范围、群体经验、军数或后果，仍须与地方文书、物质遗存及相关政权记录互证。

\paragraph{1003年}\eventanchor{NS-1003-REGNAL-YEAR}{NS-1003-REGNAL-YEAR}　宋以宋真宗在位第6年作为诏令、赋役与外交文书的纪年框架。账本列出《宋史》卷1—23〈本纪〉；《续资治通鉴长编》本年条；《宋会要辑稿》相关门类作为来源线索；该材料可核对时间与制度称谓，但有关执行范围、群体经验、军数或后果，仍须与地方文书、物质遗存及相关政权记录互证。

\paragraph{1004年冬}\eventanchor{NS-1004-CHANYUAN-BATTLE}{NS-1004-CHANYUAN-BATTLE}　辽圣宗、萧太后南下，宋真宗至澶州前线，双方议和。账本列出《宋史》卷5〈真宗本纪一〉；《辽史》〈圣宗本纪〉；《续资治通鉴长编》景德元年冬条作为来源线索；该材料可核对时间与制度称谓，但有关执行范围、群体经验、军数或后果，仍须与地方文书、物质遗存及相关政权记录互证。

\paragraph{1004年}\eventanchor{NS-1004-REGNAL-YEAR}{NS-1004-REGNAL-YEAR}　宋以宋真宗在位第7年作为诏令、赋役与外交文书的纪年框架。账本列出《宋史》卷1—23〈本纪〉；《续资治通鉴长编》本年条；《宋会要辑稿》相关门类作为来源线索；该材料可核对时间与制度称谓，但有关执行范围、群体经验、军数或后果，仍须与地方文书、物质遗存及相关政权记录互证。

\paragraph{1005年}\eventanchor{NS-1005-CHANYUAN-TREATY}{NS-1005-CHANYUAN-TREATY}　澶渊盟约落实，宋按约输送银绢，双方以兄弟之国交往。账本列出《宋史》卷5〈真宗本纪一〉；《辽史》〈圣宗本纪〉；《宋会要辑稿·蕃夷》辽国盟约条作为来源线索；该材料可核对时间与制度称谓，但有关执行范围、群体经验、军数或后果，仍须与地方文书、物质遗存及相关政权记录互证。

\paragraph{1005年}\eventanchor{NS-1005-REGNAL-YEAR}{NS-1005-REGNAL-YEAR}　宋以宋真宗在位第8年作为诏令、赋役与外交文书的纪年框架。账本列出《宋史》卷1—23〈本纪〉；《续资治通鉴长编》本年条；《宋会要辑稿》相关门类作为来源线索；该材料可核对时间与制度称谓，但有关执行范围、群体经验、军数或后果，仍须与地方文书、物质遗存及相关政权记录互证。

\paragraph{1006年}\eventanchor{NS-1006-REGNAL-YEAR}{NS-1006-REGNAL-YEAR}　宋以宋真宗在位第9年作为诏令、赋役与外交文书的纪年框架。账本列出《宋史》卷1—23〈本纪〉；《续资治通鉴长编》本年条；《宋会要辑稿》相关门类作为来源线索；该材料可核对时间与制度称谓，但有关执行范围、群体经验、军数或后果，仍须与地方文书、物质遗存及相关政权记录互证。

\paragraph{1007年}\eventanchor{NS-1007-REGNAL-YEAR}{NS-1007-REGNAL-YEAR}　宋以宋真宗在位第10年作为诏令、赋役与外交文书的纪年框架。账本列出《宋史》卷1—23〈本纪〉；《续资治通鉴长编》本年条；《宋会要辑稿》相关门类作为来源线索；该材料可核对时间与制度称谓，但有关执行范围、群体经验、军数或后果，仍须与地方文书、物质遗存及相关政权记录互证。

\paragraph{1008年}\eventanchor{NS-1008-FENGSHAN}{NS-1008-FENGSHAN}　真宗举行封禅及相关大礼。账本列出《宋史》卷6〈真宗本纪二〉；《宋会要辑稿·礼》封禅条；《续资治通鉴长编》大中祥符元年条作为来源线索；该材料可核对时间与制度称谓，但有关执行范围、群体经验、军数或后果，仍须与地方文书、物质遗存及相关政权记录互证。

\paragraph{1008年}\eventanchor{NS-1008-REGNAL-YEAR}{NS-1008-REGNAL-YEAR}　宋以宋真宗在位第11年作为诏令、赋役与外交文书的纪年框架。账本列出《宋史》卷1—23〈本纪〉；《续资治通鉴长编》本年条；《宋会要辑稿》相关门类作为来源线索；该材料可核对时间与制度称谓，但有关执行范围、群体经验、军数或后果，仍须与地方文书、物质遗存及相关政权记录互证。

\paragraph{1009年}\eventanchor{NS-1009-REGNAL-YEAR}{NS-1009-REGNAL-YEAR}　宋以宋真宗在位第12年作为诏令、赋役与外交文书的纪年框架。账本列出《宋史》卷1—23〈本纪〉；《续资治通鉴长编》本年条；《宋会要辑稿》相关门类作为来源线索；该材料可核对时间与制度称谓，但有关执行范围、群体经验、军数或后果，仍须与地方文书、物质遗存及相关政权记录互证。

\paragraph{1010年}\eventanchor{NS-1010-LIAO-INCIDENT}{NS-1010-LIAO-INCIDENT}　辽军再入宋境，宋廷强化河北边防并维持盟约渠道。账本列出《宋史》卷6〈真宗本纪二〉；《辽史》〈圣宗本纪〉；《续资治通鉴长编》大中祥符三年条作为来源线索；该材料可核对时间与制度称谓，但有关执行范围、群体经验、军数或后果，仍须与地方文书、物质遗存及相关政权记录互证。

\paragraph{1010年}\eventanchor{NS-1010-REGNAL-YEAR}{NS-1010-REGNAL-YEAR}　宋以宋真宗在位第13年作为诏令、赋役与外交文书的纪年框架。账本列出《宋史》卷1—23〈本纪〉；《续资治通鉴长编》本年条；《宋会要辑稿》相关门类作为来源线索；该材料可核对时间与制度称谓，但有关执行范围、群体经验、军数或后果，仍须与地方文书、物质遗存及相关政权记录互证。

\paragraph{1011年}\eventanchor{NS-1011-REGNAL-YEAR}{NS-1011-REGNAL-YEAR}　宋以宋真宗在位第14年作为诏令、赋役与外交文书的纪年框架。账本列出《宋史》卷1—23〈本纪〉；《续资治通鉴长编》本年条；《宋会要辑稿》相关门类作为来源线索；该材料可核对时间与制度称谓，但有关执行范围、群体经验、军数或后果，仍须与地方文书、物质遗存及相关政权记录互证。

\paragraph{1012—1013年}\eventanchor{NS-1012-CHAMPA-RICE}{NS-1012-CHAMPA-RICE}　朝廷在江淮等地推广早熟稻种，后世常称占城稻。账本列出《宋史》卷173〈食货志上一〉；《宋会要辑稿·食货》农田、稻种条；《续资治通鉴长编》大中祥符五年条作为来源线索；该材料可核对时间与制度称谓，但有关执行范围、群体经验、军数或后果，仍须与地方文书、物质遗存及相关政权记录互证。

\paragraph{1012年}\eventanchor{NS-1012-REGNAL-YEAR}{NS-1012-REGNAL-YEAR}　宋以宋真宗在位第15年作为诏令、赋役与外交文书的纪年框架。账本列出《宋史》卷1—23〈本纪〉；《续资治通鉴长编》本年条；《宋会要辑稿》相关门类作为来源线索；该材料可核对时间与制度称谓，但有关执行范围、群体经验、军数或后果，仍须与地方文书、物质遗存及相关政权记录互证。

\paragraph{1013年}\eventanchor{NS-1013-REGNAL-YEAR}{NS-1013-REGNAL-YEAR}　宋以宋真宗在位第16年作为诏令、赋役与外交文书的纪年框架。账本列出《宋史》卷1—23〈本纪〉；《续资治通鉴长编》本年条；《宋会要辑稿》相关门类作为来源线索；该材料可核对时间与制度称谓，但有关执行范围、群体经验、军数或后果，仍须与地方文书、物质遗存及相关政权记录互证。

\paragraph{1014年}\eventanchor{NS-1014-REGNAL-YEAR}{NS-1014-REGNAL-YEAR}　宋以宋真宗在位第17年作为诏令、赋役与外交文书的纪年框架。账本列出《宋史》卷1—23〈本纪〉；《续资治通鉴长编》本年条；《宋会要辑稿》相关门类作为来源线索；该材料可核对时间与制度称谓，但有关执行范围、群体经验、军数或后果，仍须与地方文书、物质遗存及相关政权记录互证。

\paragraph{1015年}\eventanchor{NS-1015-REGNAL-YEAR}{NS-1015-REGNAL-YEAR}　宋以宋真宗在位第18年作为诏令、赋役与外交文书的纪年框架。账本列出《宋史》卷1—23〈本纪〉；《续资治通鉴长编》本年条；《宋会要辑稿》相关门类作为来源线索；该材料可核对时间与制度称谓，但有关执行范围、群体经验、军数或后果，仍须与地方文书、物质遗存及相关政权记录互证。

\paragraph{1016年}\eventanchor{NS-1016-REGNAL-YEAR}{NS-1016-REGNAL-YEAR}　宋以宋真宗在位第19年作为诏令、赋役与外交文书的纪年框架。账本列出《宋史》卷1—23〈本纪〉；《续资治通鉴长编》本年条；《宋会要辑稿》相关门类作为来源线索；该材料可核对时间与制度称谓，但有关执行范围、群体经验、军数或后果，仍须与地方文书、物质遗存及相关政权记录互证。

\paragraph{1017年}\eventanchor{NS-1017-REGNAL-YEAR}{NS-1017-REGNAL-YEAR}　宋以宋真宗在位第20年作为诏令、赋役与外交文书的纪年框架。账本列出《宋史》卷1—23〈本纪〉；《续资治通鉴长编》本年条；《宋会要辑稿》相关门类作为来源线索；该材料可核对时间与制度称谓，但有关执行范围、群体经验、军数或后果，仍须与地方文书、物质遗存及相关政权记录互证。

\paragraph{1018年}\eventanchor{NS-1018-REGNAL-YEAR}{NS-1018-REGNAL-YEAR}　宋以宋真宗在位第21年作为诏令、赋役与外交文书的纪年框架。账本列出《宋史》卷1—23〈本纪〉；《续资治通鉴长编》本年条；《宋会要辑稿》相关门类作为来源线索；该材料可核对时间与制度称谓，但有关执行范围、群体经验、军数或后果，仍须与地方文书、物质遗存及相关政权记录互证。

\paragraph{1019年}\eventanchor{NS-1019-REGNAL-YEAR}{NS-1019-REGNAL-YEAR}　宋以宋真宗在位第22年作为诏令、赋役与外交文书的纪年框架。账本列出《宋史》卷1—23〈本纪〉；《续资治通鉴长编》本年条；《宋会要辑稿》相关门类作为来源线索；该材料可核对时间与制度称谓，但有关执行范围、群体经验、军数或后果，仍须与地方文书、物质遗存及相关政权记录互证。

\paragraph{1020年}\eventanchor{NS-1020-REGNAL-YEAR}{NS-1020-REGNAL-YEAR}　宋以宋真宗在位第23年作为诏令、赋役与外交文书的纪年框架。账本列出《宋史》卷1—23〈本纪〉；《续资治通鉴长编》本年条；《宋会要辑稿》相关门类作为来源线索；该材料可核对时间与制度称谓，但有关执行范围、群体经验、军数或后果，仍须与地方文书、物质遗存及相关政权记录互证。

\paragraph{1021年}\eventanchor{NS-1021-REGNAL-YEAR}{NS-1021-REGNAL-YEAR}　宋以宋真宗在位第24年作为诏令、赋役与外交文书的纪年框架。账本列出《宋史》卷1—23〈本纪〉；《续资治通鉴长编》本年条；《宋会要辑稿》相关门类作为来源线索；该材料可核对时间与制度称谓，但有关执行范围、群体经验、军数或后果，仍须与地方文书、物质遗存及相关政权记录互证。

\paragraph{1022年}\eventanchor{NS-1022-REGNAL-YEAR}{NS-1022-REGNAL-YEAR}　宋以宋真宗在位第25年作为诏令、赋役与外交文书的纪年框架。账本列出《宋史》卷1—23〈本纪〉；《续资治通鉴长编》本年条；《宋会要辑稿》相关门类作为来源线索；该材料可核对时间与制度称谓，但有关执行范围、群体经验、军数或后果，仍须与地方文书、物质遗存及相关政权记录互证。

\paragraph{1022年二月}\eventanchor{NS-1022-RENZONG-ACCESSION}{NS-1022-RENZONG-ACCESSION}　真宗去世，仁宗即位，刘太后临朝。账本列出《宋史》卷8〈仁宗本纪一〉；《宋史》卷242〈章献明肃刘皇后传〉；《续资治通鉴长编》乾兴元年条作为来源线索；该材料可核对时间与制度称谓，但有关执行范围、群体经验、军数或后果，仍须与地方文书、物质遗存及相关政权记录互证。


